% ---------------------------------------------------------------------------
% Methods: synthetic data, noise model, and the training / validation / test
% sets. Numbers checked against the data on 2026-10-12.
% Items in [square brackets] are for you to fill in; \cite keys are placeholders.
% ---------------------------------------------------------------------------

\subsection{Synthetic orientation maps}

We generated 40{,}000 two-dimensional polycrystalline microstructures with
DREAM.3D~\cite{groeber2014dream3d}, each a $128 \times 128$ pixel orientation
map on a square grid with a step of 1~$\mu$m and a single cubic phase (nickel,
point group $m\bar{3}m$) [grain-size distribution and texture settings used in
DREAM.3D]. A map contains on average 50 grains (standard deviation 4),
the mean equivalent circle diameter of a grain is about 20 pixels, and 17\% of
the pixels lie on a grain boundary, which we define as a pixel with at least
one of its four neighbours in a different grain. Two maps were incomplete
because their generation run was interrupted and were discarded, leaving
39{,}998. Since every map is a separate microstructure, we split the data by
map: maps 1--500 form the test set, maps 501--1000 the validation set, and
the remaining 38{,}998 maps the training set. No grain therefore appears in more
than one set. All maps were produced with the same generator settings, and the
number of grains per map is the same across the three ranges.

\subsection{Noise model}

Each noisy map combines two kinds of error. The first is orientation scatter:
every pixel is rotated by a small rotation about an axis drawn uniformly on the
sphere, with a rotation angle drawn uniformly between $0$ and $\theta_{\max}$.
The disorientation between a
scattered pixel and its true orientation therefore has a mean of
$\theta_{\max}/2$. The second is misindexing: a fraction $p$ of the pixels is
replaced by orientations drawn uniformly from SO(3)~\cite{shoemake1992uniform},
so that these pixels carry no information about their true orientation.
Because misindexing in real scans occurs mostly where the diffraction pattern
mixes two grains, a boundary pixel is chosen four times as often as an interior
pixel, sampled without replacement. Scatter is applied first and misindexing
second. All errors are measured as the disorientation angle, the smallest
rotation between two orientations once the 24 proper rotations of the cubic
group are taken into account.

The network is trained on noise that is drawn afresh every time a training map
is used. For each draw, $\theta_{\max}$ is sampled uniformly from
$[0^\circ, 15^\circ]$ and $p$ from $[0, 10\%]$, so the network never sees the
same corruption twice and is never told the noise level of its input. We also
apply a random rotation of the specimen frame to all orientations of the map,
which leaves every misorientation between pixels unchanged, together with
random flips and $90^\circ$ rotations of the pixel grid. The 500 validation
maps were corrupted once, with a level drawn per map from the same ranges and a
fixed random seed, so that every epoch and every training run is compared on
identical inputs. The validation set is used only to select the checkpoint.

\subsection{Test conditions}

The test maps were corrupted under 28 fixed conditions: every combination of
$\theta_{\max} \in \{1, 3, 5, 7, 9, 12, 15\}^\circ$ and
$p \in \{0, 3, 6, 9\}\%$, each applied to all 500 test maps
(Fig.~\ref{fig:noise_examples}). The lowest level gives a mean error of
$0.5^\circ$, comparable to the angular precision usually reported for
Hough-based indexing~\cite{demirel2009accuracy}; at $\theta_{\max} = 15^\circ$
(mean $7.5^\circ$) the grains are only just distinguishable. All 28
conditions were generated with the same random seed. For a given map, the
rotation axes and the fraction of $\theta_{\max}$ drawn for each pixel are thus
the same in every condition, and the scatter is identical across the four
misindexing levels, so any difference between two conditions comes from the
amount of noise and not from a different random draw. At
$\theta_{\max} \geq 12^\circ$, scatter alone moves 17--33\% of the pixels
beyond $10^\circ$ from their true orientation. We therefore assess the repair
of misindexed pixels on the recorded set of replaced pixels rather than with a
fixed angular threshold.

\begin{figure*}[t]
  \centering
  \includegraphics[width=\textwidth]{figures/noise_examples.pdf}
  \caption{The noise model on test map 1, coloured with the IPF-Z key.
    (a)~Clean map. (b)--(e)~Orientation scatter alone, with maximum angle
    $\theta_{\max} = 1^\circ$, $5^\circ$, $9^\circ$ and $15^\circ$.
    (f)--(h)~Misindexing alone, $p = 3$, $6$ and $9\%$ of the pixels: the
    replaced pixels of the corresponding test sets shown on the clean map,
    without their scatter. (i)~Positions of the misindexed pixels in (h)
    (blue) relative to the grain boundaries (grey): boundary pixels
    make up 16\% of the map but receive 41\% of the misindexing, because a
    boundary pixel is four times as likely to be chosen. (j)--(m)~Scatter and
    misindexing together, from the mildest to the most severe test condition:
    $\theta_{\max} = 1^\circ$ with $p = 3\%$, $5^\circ$ with $6\%$,
    $9^\circ$ with $9\%$, and $15^\circ$ with $9\%$.}
  \label{fig:noise_examples}
\end{figure*}
